% TOP_PROBLEM: {"id": 30006737, "title": "Serre's congruence subgroup conjecture", "category_id": 17, "category": {"id": 17, "name": "group_theory", "display_name": "Group Theory", "description": "Problems about groups, group actions, representations, and related algebraic structures.", "slug": "group-theory", "order_index": 17, "created_at": "2026-07-31T15:26:25.671Z"}, "status": "open"}
% FIRST_POSTED: 2026-09-25
% ENTRY_KIND: statement_only
% !TeX program = lualatex
% Definition and sources only; this file is not an LLM solution attempt.
\documentclass[11pt]{article}
\usepackage[margin=25mm]{geometry}
\usepackage{fontspec}
\setmainfont{Latin Modern Roman}
\usepackage{amsmath,amssymb,mathtools}
\usepackage{unicode-math}
\setmathfont{Latin Modern Math}
\usepackage{xurl,hyperref}
\hypersetup{colorlinks=true,urlcolor=blue}
\setcounter{secnumdepth}{0}
\setlength{\parindent}{0pt}
\setlength{\parskip}{5pt}
\setlength{\emergencystretch}{3em}
\begin{document}
\section*{\texorpdfstring{Serre's congruence subgroup conjecture}{Serre's congruence subgroup conjecture}}\label{30006737}
\subsection{Definitions and mathematical statement}
For the global field $K$ and finite place set $S$, let $\mathcal O_S$ be the ring of elements integral outside $S$. The congruence completion of $G(\mathcal O_S)$ uses kernels of reduction modulo ideals, whereas its profinite completion uses all finite-index normal subgroups. Define
\[
 C^S(G)=\ker\bigl(\widehat{G(\mathcal O_S)}\to\overline{G(\mathcal O_S)}^{\rm cong}\bigr),
 \qquad r_S(G)=\sum_{v\in S}\operatorname{rank}_{K_v}G.
\]
For absolutely almost simple simply connected $G$, assert finiteness of this kernel when $r_S\ge2$ and every nonarchimedean place in $S$ has positive local rank; assert infinitude when $r_S=1$. Rank zero is excluded. Finiteness, not triviality, is the higher-rank conclusion.
\par\smallskip\textit{Formulation and concept references:} \href{https://www.mathi.uni-heidelberg.de/~wienhard/retreat10/references/prasad_rapinchuk_CSP.pdf}{[S1]}.

\subsection{Short English statement}
Does the congruence kernel have the finite-versus-infinite behavior predicted by the total local rank of an arithmetic group?
\subsection{Sources}
\begin{itemize}
\item[S1]G. Prasad and A. S. Rapinchuk, Developments on the congruence subgroup problem after the work of Bass, Milnor and Serre.
\url{https://www.mathi.uni-heidelberg.de/~wienhard/retreat10/references/prasad_rapinchuk_CSP.pdf}.
\textit{Formulation source.}
\end{itemize}
\end{document}
